\chapter{Tiền Đầu Tiên, Tự Do Đầu Tiên, Sai Lầm Đầu Tiên}
\markboth{Tiền Đầu Tiên, Tự Do Đầu Tiên, Sai Lầm Đầu Tiên}{First Money, First Freedom, First Mistakes}

\section{Lương Đầu Không Nhiều Như Em Tưởng}
\begin{truyen}{Lương Đầu Không Nhiều Như Em Tưởng}{The First Salary Feels Smaller Than You Thought}
\chuhoa{N}{hiều người tưởng cầm lương đầu sẽ thấy mình lớn hẳn lên.}
\textit{(Many people think their first salary will make them feel fully grown.)}

Nhưng vừa trừ tiền nhà, tiền ăn, đi lại, và vài khoản phát sinh, số còn lại rất khiêm tốn.
\textit{(But after rent, food, transport, and surprise expenses, what remains is often modest.)}

Lúc đó nhiều người mới hiểu vì sao người lớn quanh mình luôn phải tính.
\textit{(At that moment many finally understand why the adults around them always had to calculate carefully.)}

Tiền đầu tiên không chỉ cho em niềm vui.
\textit{(The first salary gives not only joy.)}

Nó cho em cái nhìn thật hơn về đời sống.
\textit{(It gives a more realistic view of life.)}

\end{truyen}

\begin{baihoc}
Kiếm được tiền là một bước lớn. Nhưng học cách giữ tiền mới là bước khiến em bớt ngây thơ.
\textit{(Earning money is a big step. Learning how to keep it is the step that makes you less naive.)}
\end{baihoc}

\begin{ghinhoanh}
\textbf{Short English to remember:}

Earning money is a big step.

Learning how to keep it is the step that makes you less naive.

\end{ghinhoanh}

\begin{tuhocnhanh}
\textbf{Cau hoi tu hoc / Self-study questions}

1. Van de chinh la gi? \textit{What is the main problem?}

2. Bai hoc thuc te la gi? \textit{What is the practical lesson?}

3. Mot cau tieng Anh em muon nho la cau nao? \textit{Which English sentence do you want to remember?}
\end{tuhocnhanh}
\ngancach

\section{Tiền Tự Do Và Cái Bẫy Mua Để Chứng Minh}
\begin{truyen}{Tiền Tự Do Và Cái Bẫy Mua Để Chứng Minh}{Freedom Money and the Trap of Buying to Prove Yourself}
\chuhoa{K}{hi có tiền đầu tiên, nhiều người muốn mua một thứ gì đó để thấy mình không còn là đứa trẻ đi xin nữa.}
\textit{(When people get their first money, many want to buy something just to feel they are no longer children asking for things.)}

Mong muốn đó rất dễ hiểu.
\textit{(That desire is easy to understand.)}

Nhưng nếu không tỉnh, em sẽ tiêu để chứng minh hơn là tiêu để sống.
\textit{(But without awareness, you spend to prove something more than to live well.)}

Một chiếc điện thoại đắt có thể cho em cảm giác hơn người trong vài ngày.
\textit{(An expensive phone may give you a sense of status for a few days.)}

Một khoản tiết kiệm lại cho em bình an nhiều tháng.
\textit{(Savings can give peace for many months.)}

Người trẻ hay lẫn giữa cái mình thích ngay và cái mình cần lâu dài.
\textit{(Young people often confuse what they want now with what they need for the long run.)}

\end{truyen}

\begin{baihoc}
Tiền đầu tiên rất dễ bị tiêu vì cái tôi. Muốn trưởng thành tài chính, em phải học tiêu theo giá trị, không chỉ theo cảm xúc.
\textit{(First money is easily spent by the ego. To grow financially, you must spend by values, not only by emotion.)}
\end{baihoc}

\begin{ghinhoanh}
\textbf{Short English to remember:}

First money is easily spent by the ego.

To grow financially, you must spend by values, not only by emotion.

\end{ghinhoanh}

\begin{tuhocnhanh}
\textbf{Cau hoi tu hoc / Self-study questions}

1. Van de chinh la gi? \textit{What is the main problem?}

2. Bai hoc thuc te la gi? \textit{What is the practical lesson?}

3. Mot cau tieng Anh em muon nho la cau nao? \textit{Which English sentence do you want to remember?}
\end{tuhocnhanh}
\ngancach

\section{Một Lần Quẹt Sai, Mất Cả Tháng Thở Dài}
\begin{truyen}{Một Lần Quẹt Sai, Mất Cả Tháng Thở Dài}{One Bad Swipe Can Cost You a Whole Month}
\chuhoa{C}{ó người trẻ dùng thẻ hay vay app rất nhẹ tay vì tiền chưa phải rời khỏi ví ngay trước mắt.}
\textit{(Some young people use cards or loan apps too easily because the money does not seem to leave their hand right away.)}

Đến lúc phải trả, họ mới thấy quyết định vài giây có thể kéo dài áp lực cả tháng.
\textit{(When payment time arrives, they realize that a few seconds of decision can create a month of pressure.)}

Tiện lợi tài chính là thứ đẹp ở mặt ngoài nhưng dễ cắn nếu dùng thiếu tỉnh táo.
\textit{(Financial convenience looks nice on the surface but bites hard when used without care.)}

Nhiều người lớn khổ không phải vì họ không làm ra tiền. Họ khổ vì không quản được nhịp tiêu tiền.
\textit{(Many adults suffer not because they cannot earn money, but because they cannot manage the rhythm of spending.)}

\end{truyen}

\begin{baihoc}
Đừng coi nợ nhỏ là vô hại. Nhiều cái vòng siết bắt đầu từ những khoản rất nhỏ nhưng rất dễ dãi.
\textit{(Do not treat small debt as harmless. Many traps begin with tiny amounts handled too casually.)}
\end{baihoc}

\begin{ghinhoanh}
\textbf{Short English to remember:}

Do not treat small debt as harmless.

Many traps begin with tiny amounts handled too casually.

\end{ghinhoanh}

\begin{tuhocnhanh}
\textbf{Cau hoi tu hoc / Self-study questions}

1. Van de chinh la gi? \textit{What is the main problem?}

2. Bai hoc thuc te la gi? \textit{What is the practical lesson?}

3. Mot cau tieng Anh em muon nho la cau nao? \textit{Which English sentence do you want to remember?}
\end{tuhocnhanh}
\ngancach

\section{Tiền Đầu Không Chỉ Là Tiền Của Em}
\begin{truyen}{Tiền Đầu Không Chỉ Là Tiền Của Em}{Your First Money Is Not Only About You}
\chuhoa{N}{hiều người trẻ lần đầu cầm tiền mới thấy một cảm giác lạ: mình muốn mua cho bố mẹ cái gì đó.}
\textit{(Many young people feel something strange with their first money: they want to buy something for their parents.)}

Khoảnh khắc đó đánh dấu một bước tâm lý rất lớn.
\textit{(That moment marks a major psychological step.)}

Em bắt đầu hiểu mình không còn chỉ là người nhận.
\textit{(You begin to understand that you are no longer only a receiver.)}

Trưởng thành một phần là từ chỗ nghĩ mình cần gì sang chỗ nghĩ mình có thể đỡ ai được gì.
\textit{(Part of growing up is moving from what you need to what you can help carry for others.)}

Tiền đầu tiên vì thế không chỉ có giá trị vật chất. Nó còn có giá trị đạo đức.
\textit{(That is why first money holds not only material value, but moral value too.)}

\end{truyen}

\begin{baihoc}
Khi đồng tiền đầu tiên khiến em nghĩ đến người khác, đó là dấu hiệu em đang lớn lên đúng hướng.
\textit{(When your first money makes you think about others, it is a sign that you are growing in the right direction.)}
\end{baihoc}

\begin{ghinhoanh}
\textbf{Short English to remember:}

When your first money makes you think about others, it is a sign that you are growing in the right direction.

\end{ghinhoanh}

\begin{tuhocnhanh}
\textbf{Cau hoi tu hoc / Self-study questions}

1. Van de chinh la gi? \textit{What is the main problem?}

2. Bai hoc thuc te la gi? \textit{What is the practical lesson?}

3. Mot cau tieng Anh em muon nho la cau nao? \textit{Which English sentence do you want to remember?}
\end{tuhocnhanh}
\ngancach

\section{Tự Do Không Có Nghĩa Là Muốn Tiêu Sao Cũng Được}
\begin{truyen}{Tự Do Không Có Nghĩa Là Muốn Tiêu Sao Cũng Được}{Freedom Does Not Mean Spending Without Limits}
\chuhoa{N}{hiều người trẻ lẫn giữa tự do và tùy tiện.}
\textit{(Many young people confuse freedom with impulsiveness.)}

Có tiền riêng rồi, họ nghĩ mình có quyền sống theo cảm xúc hoàn toàn.
\textit{(Once they earn their own money, they think they have the right to live entirely by feeling.)}

Nhưng tự do thật không phá hỏng tương lai của mình chỉ để thỏa mãn hiện tại.
\textit{(But real freedom does not damage your future just to satisfy the present.)}

Người không kiểm soát được tiêu dùng thường sớm bị mất tự do bởi chính những lựa chọn của mình.
\textit{(People who cannot control spending often lose freedom to their own choices very quickly.)}

\end{truyen}

\begin{baihoc}
Tiền có thể cho em lựa chọn. Nhưng chỉ kỷ luật mới giữ cho những lựa chọn đó không quay lại bóp nghẹt em.
\textit{(Money can give you options. But only discipline keeps those options from turning around and choking you.)}
\end{baihoc}

\begin{ghinhoanh}
\textbf{Short English to remember:}

Money can give you options.

But only discipline keeps those options from turning around and choking you.

\end{ghinhoanh}

\begin{tuhocnhanh}
\textbf{Cau hoi tu hoc / Self-study questions}

1. Van de chinh la gi? \textit{What is the main problem?}

2. Bai hoc thuc te la gi? \textit{What is the practical lesson?}

3. Mot cau tieng Anh em muon nho la cau nao? \textit{Which English sentence do you want to remember?}
\end{tuhocnhanh}
\ngancach

\section{Quỹ Lương Đầu Cho Nhà}
\begin{truyen}{Quỹ Lương Đầu Cho Nhà}{Giving Part of the First Salary to Home}
\chuhoa{C}{ó người cầm lương đầu rồi nghĩ ngay đến việc phụ bố mẹ một khoản nhỏ.}
\textit{(Some people receive their first salary and immediately think about helping their parents with a small part.)}

Trước đó họ tưởng tiền đầu chỉ là phần thưởng cho riêng mình.
\textit{(Before that, they thought the first salary was only a reward for themselves.)}

Nhưng khi bắt đầu nghĩ đến gia đình, họ cảm thấy một tầng trưởng thành khác mở ra.
\textit{(But when they begin to think about family, another layer of maturity opens.)}

Đồng tiền đầu tiên không chỉ thay đổi ví tiền.
\textit{(First money changes not only the wallet.)}

Nó đổi cách mình nhìn vai trò của mình trong nhà.
\textit{(It changes how you see your role at home.)}

\end{truyen}

\begin{baihoc}
Lúc nghĩ đến việc đỡ người khác bằng đồng tiền mình làm ra, em bắt đầu lớn thật hơn.
\textit{(When your earned money starts serving others too, you grow more truly.)}
\end{baihoc}

\begin{ghinhoanh}
\textbf{Short English to remember:}

When your earned money starts serving others too, you grow more truly.

\end{ghinhoanh}

\begin{tuhocnhanh}
\textbf{Cau hoi tu hoc / Self-study questions}

1. Van de chinh la gi? \textit{What is the main problem?}

2. Bai hoc thuc te la gi? \textit{What is the practical lesson?}

3. Mot cau tieng Anh em muon nho la cau nao? \textit{Which English sentence do you want to remember?}
\end{tuhocnhanh}
\ngancach

\section{Quỹ Dự Phòng Đầu Tiên}
\begin{truyen}{Quỹ Dự Phòng Đầu Tiên}{The First Emergency Fund}
\chuhoa{S}{au vài cú hụt tiền, có người bắt đầu để riêng một khoản rất nhỏ và gọi đó là quỹ dự phòng.}
\textit{(After a few money shortages, someone starts keeping a small amount aside as an emergency fund.)}

Họ từng tưởng tiết kiệm là chuyện của người lớn tuổi hoặc người nhiều tiền.
\textit{(They once thought saving was for older people or for those with more money.)}

Nhưng chính người chưa có nhiều lại càng cần một khoảng đệm nhỏ để bớt hoảng khi có việc.
\textit{(But those with less often need a small buffer even more to avoid panic when problems come.)}

Một khoản nhỏ để yên có thể cứu em khỏi rất nhiều quyết định vội vàng.
\textit{(A small amount left untouched can save you from many rushed decisions.)}

\end{truyen}

\begin{baihoc}
Tiền dự phòng không làm em già đi. Nó làm em bớt ngây thơ.
\textit{(An emergency fund does not make you old. It makes you less naive.)}
\end{baihoc}

\begin{ghinhoanh}
\textbf{Short English to remember:}

An emergency fund does not make you old.

It makes you less naive.

\end{ghinhoanh}

\begin{tuhocnhanh}
\textbf{Cau hoi tu hoc / Self-study questions}

1. Van de chinh la gi? \textit{What is the main problem?}

2. Bai hoc thuc te la gi? \textit{What is the practical lesson?}

3. Mot cau tieng Anh em muon nho la cau nao? \textit{Which English sentence do you want to remember?}
\end{tuhocnhanh}
\ngancach

\section{Mua Để Bằng Bạn Bằng Bè}
\begin{truyen}{Mua Để Bằng Bạn Bằng Bè}{Buying Things to Keep Up with Friends}
\chuhoa{T}{hấy bạn bè có đồ mới, có quán mới, có cuộc vui mới, nhiều người trẻ rất khó đứng yên.}
\textit{(Seeing friends with new things, new places, and new outings makes many young people restless.)}

Họ tưởng theo kịp hình ảnh của nhóm là một phần của tự do tài chính.
\textit{(They think keeping up with the group image is part of financial freedom.)}

Nhưng tiêu để bằng người khác thường chỉ mua được vài giờ dễ chịu rồi đổi lấy nhiều ngày áp lực.
\textit{(But spending to match others often buys a few easy hours and many stressed days later.)}

Người biết giữ tiền thường chịu được cảm giác không bằng ai đó trong một lúc.
\textit{(Those who keep money well can tolerate not matching someone else for a while.)}

\end{truyen}

\begin{baihoc}
Trưởng thành tài chính bắt đầu từ lúc em ngừng tiêu để chứng minh mình không thua.
\textit{(Financial maturity begins when you stop spending to prove you are not behind.)}
\end{baihoc}

\begin{ghinhoanh}
\textbf{Short English to remember:}

Financial maturity begins when you stop spending to prove you are not behind.

\end{ghinhoanh}

\begin{tuhocnhanh}
\textbf{Cau hoi tu hoc / Self-study questions}

1. Van de chinh la gi? \textit{What is the main problem?}

2. Bai hoc thuc te la gi? \textit{What is the practical lesson?}

3. Mot cau tieng Anh em muon nho la cau nao? \textit{Which English sentence do you want to remember?}
\end{tuhocnhanh}
\ngancach

\section{Giá Của Sống Riêng}
\begin{truyen}{Giá Của Sống Riêng}{The Price of Living Independently}
\chuhoa{C}{ó người chỉ khi tự trả tiền nhà, tiền ăn, và tiền hỏng vặt mới hiểu vì sao người lớn luôn cân đo.}
\textit{(Only when paying rent, food, and small repair costs do some people understand why adults always calculate.)}

Trước đó họ tưởng sống riêng chủ yếu là cảm giác thoải mái và ít bị quản.
\textit{(Before that, they thought living independently was mostly comfort and less control.)}

Nhưng mỗi sự tiện nghi đều có giá thật, và giá đó lặp lại từng tháng.
\textit{(But every convenience has a real cost, and that cost returns each month.)}

Tự do nhìn rất đẹp cho đến khi em phải tự trả toàn bộ hóa đơn của nó.
\textit{(Freedom looks beautiful until you pay its full bill yourself.)}

\end{truyen}

\begin{baihoc}
Biết giá của cuộc sống là bước đầu để biết quý đồng tiền mình làm ra.
\textit{(Knowing the cost of life is the first step to valuing your earned money.)}
\end{baihoc}

\begin{ghinhoanh}
\textbf{Short English to remember:}

Knowing the cost of life is the first step to valuing your earned money.

\end{ghinhoanh}

\begin{tuhocnhanh}
\textbf{Cau hoi tu hoc / Self-study questions}

1. Van de chinh la gi? \textit{What is the main problem?}

2. Bai hoc thuc te la gi? \textit{What is the practical lesson?}

3. Mot cau tieng Anh em muon nho la cau nao? \textit{Which English sentence do you want to remember?}
\end{tuhocnhanh}
\ngancach

\section{Gửi Tiền Về Nhà}
\begin{truyen}{Gửi Tiền Về Nhà}{Sending Money Home}
\chuhoa{M}{ột số người trẻ chỉ gửi về một khoản rất nhỏ, nhưng cảm giác bên trong lại lớn hơn nhiều.}
\textit{(Some young people send home only a small amount, yet the feeling inside is much larger.)}

Họ từng tưởng trưởng thành là cầm tiền cho mình và mua điều mình thích.
\textit{(They once thought adulthood meant holding money for themselves and buying what they wanted.)}

Đến lúc chủ động san sẻ, họ mới thấy đồng tiền còn mang cả ý nghĩa biết ơn.
\textit{(When they start sharing, they see that money also carries gratitude.)}

Sự trưởng thành đẹp không chỉ nằm ở kiếm được bao nhiêu, mà còn ở mình nhớ đến ai khi kiếm được.
\textit{(Beautiful maturity is not only about how much you earn, but whom you remember when you earn.)}

\end{truyen}

\begin{baihoc}
Kiếm tiền cho mình là bước đầu. Biết chia đồng tiền đó đúng chỗ là bước sâu hơn.
\textit{(Earning for yourself is the first step. Knowing how to share it wisely is a deeper one.)}
\end{baihoc}

\begin{ghinhoanh}
\textbf{Short English to remember:}

Earning for yourself is the first step.

Knowing how to share it wisely is a deeper one.

\end{ghinhoanh}

\begin{tuhocnhanh}
\textbf{Cau hoi tu hoc / Self-study questions}

1. Van de chinh la gi? \textit{What is the main problem?}

2. Bai hoc thuc te la gi? \textit{What is the practical lesson?}

3. Mot cau tieng Anh em muon nho la cau nao? \textit{Which English sentence do you want to remember?}
\end{tuhocnhanh}
